\chapter[Band 11: Äquivalenzrelationen]{Band 11 auf einen Blick}
\noindent\textbf{Äquivalenzrelationen}\par\medskip
\noindent\emph{Lesart.} Für Klassen- und Quotientenformeln gilt
\(\EqRel{A,\sim}\). Es werden Relationsoperatoren auf einem ausdrücklich
angegebenen Träger benutzt. Gleichmächtigkeit wird auf Mengenfamilien
als Beispiel einer Äquivalenzrelation betrachtet.
\begin{BandUebersicht}
\textbf{Äquivalenzrelation}
Für \(x,y,z\in A\) lauten die Strukturaxiome:
\UebFormel{\begin{aligned}
&x\sim x,\\
&x\sim y\vdash y\sim x,\\
&x\sim y\dsep y\sim z\vdash x\sim z.
\end{aligned}}
\UebQuellen{\FormulaRefAuto{EqRelDef}[def];
\FormulaRefAuto{EqRelReflexivity}[ax];
\FormulaRefAuto{EqRelSymmetry}[ax];
\FormulaRefAuto{EqRelTransitivity}[ax]}
&
\textit{Ein konkretes Modell: Ununterscheidbarkeit}
Die in der nächsten Zeile definierte Relation erfüllt alle drei Axiome:
\UebFormel{\EqRel{\bigcup M,\SameOccRel{M}}.}
\UebQuellen{\FormulaRefAuto{SameOccRelEqRel}[thm]}\\
\addlinespace[.8ex]
\textbf{Ununterscheidbarkeit in einer Familie}
Für \(a,b\in\bigcup M\):
\UebFormel{\begin{aligned}
&a\mathrel{\SameOccRel{M}}b\coloneqq\\
&\qquad\forall X\in M\,(a\in X\leftrightarrow b\in X).
\end{aligned}}
\UebQuellen{\FormulaRefAuto{SameOccRelDef}[def]}
&
\textit{Die Quotientenabbildung trennt Familienglieder}
Mit der später eingeführten Quotientenbildabbildung:
\UebFormel{\begin{aligned}
&X,Y\in M\dsep
\OccQuotMap{M}(X)=\OccQuotMap{M}(Y)\\
&\hspace{12mm}\vdash X=Y.
\end{aligned}}
\UebQuellen{\FormulaRefAuto{OccQuotSetSeparatesMembers}[thm]}\\
\addlinespace[.8ex]
\textbf{Gleichmächtigkeit}
\UebFormel{A\EqCard B\coloneqq\exists F\colon A\bij B.}
\UebQuellen{\FormulaRefAuto{Gleichmächtigkeit}[def]}
&
\textit{Äquivalenzgesetze und Potenzmengentransport}
\UebFormel{\begin{aligned}
&A\EqCard A,\qquad A\EqCard B\vdash B\EqCard A,\\
&A\EqCard B\dsep B\EqCard C\vdash A\EqCard C,\\
&A\EqCard B\vdash\powerset(A)\EqCard\powerset(B).
\end{aligned}}
\UebQuellen{\FormulaRefAuto{EqCardReflexive}[thm];
\FormulaRefAuto{EqCardSymmetric}[thm];
\FormulaRefAuto{EqCardTransitive}[thm];
\FormulaRefAuto{EqCardPowerSetForward}[thm]}\\
\addlinespace[.8ex]
\textbf{Kardinaler Größenvergleich}
\UebFormel{\begin{aligned}
A\CardLeq B&\coloneqq\exists F\colon A\inj B,\\
A\CardLt B&\coloneqq
(A\CardLeq B\land\neg(A\EqCard B)).
\end{aligned}}
\UebQuellen{\FormulaRefAuto{CardLeqDef}[def];
\FormulaRefAuto{CardLtDef}[def]}
&
\textit{Cantor--Schröder--Bernstein}
\UebFormel{\begin{aligned}
&A\CardLeq B,\ B\CardLeq A\vdash A\EqCard B,\\
&A\EqCard B\leftrightarrow
(A\CardLeq B\land B\CardLeq A).
\end{aligned}}
Der Vergleich ist reflexiv und transitiv; sein strikter Teil
ist irreflexiv.
\UebQuellen{\FormulaRefAuto{CantorSchroederBernstein}[thm];
\FormulaRefAuto{EqCardMutualCardLeqEquiv}[thm];
\FormulaRefAuto{CardLeqReflexive}[thm];
\FormulaRefAuto{CardLeqTransitive}[thm];
\FormulaRefAuto{CardLtIrreflexive}[thm]}
\\ \addlinespace[.8ex]

\textbf{Äquivalenzklasse}
Für \(x\in A\):
\UebFormel{\EqClass{x}{\sim}{A}\coloneqq
\{y\in A\mid x\sim y\}.}
\UebQuellen{\FormulaRefAuto{EqClassDef}[def]}
&
\textit{Repräsentanten und gleiche Klassen}
Für \(x,y\in A\):
\UebFormel{\begin{aligned}
&x\in\EqClass{x}{\sim}{A},\\
&\EqClass{x}{\sim}{A}=\EqClass{y}{\sim}{A}
\leftrightarrow x\sim y.
\end{aligned}}
\UebQuellen{\FormulaRefAuto{EqClassRepresentative}[thm];
\FormulaRefAuto{EqClassEqualIffRel}[thm]}\\
\addlinespace[.8ex]
\textbf{Quotientenmenge und Projektion}
\UebFormel{\begin{aligned}
\QuotientSet{A}{\sim}\coloneqq{}&
\{C\in\powerset(A)\mid{}\\
&\exists x\in A\,C=\EqClass{x}{\sim}{A}\},\\
\QuotProj{A}{\sim}(x)\coloneqq{}&\EqClass{x}{\sim}{A}.
\end{aligned}}
\UebQuellen{\FormulaRefAuto{QuotientSetDef}[def];
\FormulaRefAuto{QuotProjDef}[def]}
&
\textit{Surjektivität und identifizierte Punkte}
\UebFormel{\begin{aligned}
&\QuotProj{A}{\sim}\colon A\sur\QuotientSet{A}{\sim},\\
&x,y\in A\vdash\\
&\quad\QuotProj{A}{\sim}(x)=\QuotProj{A}{\sim}(y)
\leftrightarrow x\sim y.
\end{aligned}}
\UebQuellen{\FormulaRefAuto{QuotProjSurjective}[thm];
\FormulaRefAuto{QuotProjEqualIffRel}[thm]}\\
\addlinespace[.8ex]
\textbf{Quotientenbild einer Mengenfamilie}
Für \(M\subseteq\powerset(A)\):
\UebFormel{\QuotFamMap{A}{\sim}{M}\coloneqq
\PowMap{\QuotProj{A}{\sim}}\restriction_M.}
\UebQuellen{\FormulaRefAuto{QuotFamMapDef}[def]}
&
\textit{Vereinigungen werden erhalten}
Für \(X,Y,X\cup Y\in M\):
\UebFormel{\begin{aligned}
&\QuotFamMap{A}{\sim}{M}(X\cup Y)\\
&\quad=\QuotFamMap{A}{\sim}{M}(X)
\cup\QuotFamMap{A}{\sim}{M}(Y).
\end{aligned}}
\UebQuellen{\FormulaRefAuto{QuotFamMapPreservesUnion}[thm]}\\
\addlinespace[.8ex]
\textbf{Ununterscheidbarkeitsquotient}
\(\OccQuotCarrier{M}\) ist der Träger der Ununterscheidbarkeitsklassen,
\(\OccQuotFam{M}\) die Familie der Bilder der Glieder von \(M\).
\(\OccQuotMap{M}\) ist die zugehörige Quotientenbildabbildung.
&
\textit{Bijektive Reduktion der Mengenfamilie}
\UebFormel{\begin{aligned}
&\OccQuotMap{M}\colon M\bij\OccQuotFam{M},\\
&\UnionClosedFam(M)\\
&\hspace{10mm}\vdash\UnionClosedFam(\OccQuotFam{M}),\\
&\TraceMap{\OccQuotMap{M}}\colon
\OccQuotCarrier{M}\inj\powerset(M).
\end{aligned}}
\UebQuellen{\FormulaRefAuto{OccQuotMapBijective}[thm];
\FormulaRefAuto{OccQuotUnionClosed}[thm];
\FormulaRefAuto{OccQuotTraceMapInjective}[thm]}\\
\end{BandUebersicht}
\noindent\emph{Lesepfad zur Gleichmächtigkeit.}
Die \CSBReadingLink{} erklärt, wie aus Injektionen in beide Richtungen
eine Bijektion entsteht und damit Gleichmächtigkeit folgt;
die \CSBProofsLink{} dokumentieren die zugrunde liegende Konstruktion.
